\vspace{-8pt}
\section{Limitations and Broader Impact}
\label{sec:limitations}
% Our experiments use controlled Reddit-style simulations, enabling paired counterfactual evaluation but not capturing all deployment complexity, such as longer histories, tools, retrieval, or persistent memory. We therefore view the results as evidence of a mechanism and mitigation principle---sanitize state before summarizing it---rather than as a deployment benchmark. Our measurements rely on Detoxify, so SPG and P95 toxicity should be validated with other classifiers, human judgments, and task-specific harm measures. Finally, DPO is evaluated on \texttt{Llama-3.1-8B-Instruct} while the strongest phenomenon results use \texttt{gpt-4o-mini}, reflecting the need for parameter access; future work should replicate on more open-weight models and deployed memory architectures.
Our controlled Reddit-style simulations enable paired counterfactual evaluation but omit deployment complexity (longer histories, tools, retrieval, persistent memory); we therefore present the results as evidence of a mechanism and mitigation principle---sanitize before summarizing---rather than a deployment benchmark. Toxicity measurements rely on Detoxify and should be validated with other classifiers, human judgments, and task-specific harm measures. DPO requires parameter access and is evaluated only on \texttt{Llama-3.1-8B-Instruct}; future work should replicate on more open-weight models and deployed memory architectures.

The broader impact is primarily defensive: memory-augmented agents should be evaluated not only by outputs, but also by the summaries and state future agents consume. A dual-use risk is that memory laundering could inform attacks that hide harmful influence in benign-looking summaries; we reduce this risk by emphasizing aggregate metrics and mitigations rather than deployable attack procedures.
